
\subsubsection{Findings}

\textbf{Finding 1: Pre-alignment confidence margin is a statistically robust predictor of argmax flip for DPO-aligned tulu-2-7b.} The margin-flip relationship ($\beta _1 = -4.33, \eta^2$ = 0.289, AUROC = 0.867--0.909) is large in both statistical and practical terms for pair 2. The negative direction --- lower margin predicts higher flip probability --- is consistent with the geometric interpretation that low-margin items lie near the argmax decision boundary. Cross-benchmark generalization (predictor trained on MMLU, evaluated on TruthfulQA and ARC-Challenge without retraining) holds across all three benchmarks above the 0.75 AUROC threshold. However, the evidence for this finding is concentrated in a single model pair (tulu-2-7b DPO); pair 4 (pythia-6.9b SFT) shows a directionally consistent but substantially weaker result (AUROC = 0.609, $\eta^2$ = 0.028) that does not clear the 0.75 threshold.

\textbf{Finding 2: Alignment-induced logit perturbations are non-isotropic under both DPO and SFT.} The anisotropy ratios (pair 2 DPO: 2.90; pair 4 SFT: 4.58) are consistently and substantially above the isotropic Gaussian control (1.13) across all three benchmarks and both model pairs. This indicates that alignment-induced changes to the 4D logit space are not well-approximated by isotropic noise. The principal axis of perturbation accounts for a larger share of total perturbation variance than the remaining three axes combined in the SFT case. This finding is consistent with the geometric interpretation of the margin predictor, as structured perturbations along preferred axes would disproportionately affect items near the boundary along those axes.

\textbf{Finding 3: The H-M2 directional hypothesis is not supported; a contrasting quintile profile is observed.} DPO does not concentrate perturbation variance in low-confidence regions. The observed DPO quintile profile (monotone increase from Q1 to Q5) is opposite to the predicted direction. SFT produces a flat profile. This null result for the directional claim is documented as a limitation. The contrasting profile pattern is reported as a descriptive observation; its mechanistic interpretation and generalizability require further investigation, including a controlled experiment using the same base model aligned with both DPO and SFT.

\subsubsection{Limitations}

\textbf{L1: PPO models unavailable.} The original experimental design called for comparison between DPO, SFT, and PPO. The two planned PPO model pairs were unavailable (tulu-2-ppo-7b: HTTP 404; reciprocate/ppo\_hh\_pythia-1B: tokenizer error). All method comparisons in this paper are between DPO and SFT rather than between DPO and PPO. Conclusions about method-specific behavior cannot be extended to RLHF-based methods in general.

\textbf{L2: Primary evidence from one model pair.} The strong margin-flip result ($\beta _1 = -4.33$, AUROC = 0.867--0.909, $\eta^2$ = 0.289) derives from pair 2 alone (tulu-2-7b DPO). Pair 4 (pythia-6.9b SFT) replicates the directional effect but not the discriminative strength (AUROC = 0.609). The two pairs differ in alignment procedure, architecture, and training data; isolating the contribution of alignment procedure requires a same-base-model comparison. Claims about the generality of the margin-flip relationship beyond this specific model pair should be treated as preliminary.

\textbf{L3: H-M3 and H-M4 not executed.} The planned interaction-term analysis (H-M3: $margin\times method$ interaction in mixed-effects logistic regression) and cosine projection analysis (H-M4: alignment of logit-delta direction with the base model decision axis) were not run. The mechanistic account of why the margin predictor works therefore remains partially inferential and cannot be confirmed from the available data.

\textbf{L4: H-M2 directional hypothesis not supported.} The predicted mechanism --- that DPO concentrates perturbation energy in low-confidence regions --- is not supported in any of the three datasets. The observed pattern (DPO higher variance at high confidence) may reflect model identity confounds or a different underlying mechanism.

\textbf{L5: MCQ-specific scope.} All experiments are conducted on 4-option MCQ benchmarks. Only items producing valid 4-token logit vectors are retained. The margin predictor as operationalized here does not directly apply to free-form generation, where the output space is the full vocabulary and the concept of argmax flip across a fixed option set is not defined. Generalization to free-form settings requires a separate operationalization and empirical validation.

\textbf{L6: Model identity confounds in quintile comparison.} The quintile variance profiles are compared between pair 2 (tulu-2-7b) and pair 4 (pythia-6.9b), which differ in base model, alignment procedure, and training data. Any apparent difference in quintile profile may reflect architecture or data differences rather than the DPO vs. SFT distinction.

\subsubsection{Practical Implications}

If the margin-flip relationship generalizes beyond pair 2, the pre-alignment confidence margin could serve as a basis for item-level pre-deployment risk assessment: before deploying an aligned model, practitioners could compute base model margins for evaluation items and flag low-margin items for targeted post-alignment review. The flip rate in the bottom confidence quintile of pair 2 is approximately 25\%, compared to approximately 1.5\% in the top quintile. Whether this translates to a reliable auditing tool depends on whether the margin-flip relationship holds across additional model pairs, alignment procedures, and evaluation domains --- questions that the present data do not fully resolve.

